\begin{textsample}{POS Dim 2 – 2023\_c6 – Score 2.00 – 2023\_c6/enade\_2023\_zoot\_a5e91995a2.txt}  \label{ex:f2_pos_167}
CAPÍTULO 5 ANÁLISE TÉCNICA DA PROVA Este capítulo tem o objetivo de apresentar o desempenho dos estudantes concluintes de Zootecnia no Enade 2023.
Para isso, foram calculadas as estatísticas básicas da prova em seu todo ( subseção 5.1.1 ), bem como as estatísticas relacionadas aos componentes de Formação Geral ( subseção 5.1.2 ) e de Conhecimento Específico da área ( subseção 5.1.3 ).
Dadas as suas características, foram analisadas, em separado, as questões objetivas ( seção 5.2 ) e as questões discursivas ( seção 5.3 ).
Tomando -se como base a questão discursiva do componente de Formação Geral, nas subseções de 5.3.1.4 a 5.3.1.6, são apresentados os resultados e os comentários sobre a correção das respostas em relação ao desempenho linguístico.
Para as questões objetivas, foram disponibilizados os índices de facilidade e de discriminação ponto-bisserial, também separadamente, para os componentes de Formação Geral ( subseção 5.2.1 ) e de Conhecimento Específico ( subseção 5.2.2 ).
De cada componente, uma das questões foi escolhida para exemplificar a análise gráfica, relacionando -se as alternativas escolhidas pelos estudantes ( inclusive o gabarito ) ao número de acertos no componente.
No Anexo IV, apresenta -se a íntegra da análise gráfica para todas as questões objetivas.
Para cada uma das questões discursivas, os conteúdos dos tipos mais comuns de respostas dos estudantes são apresentados e comparados com o padrão de resposta esperado ( ver Anexo VIII ).
Nas tabelas, constam as seguintes estatísticas das notas⁵ : média do desempenho na prova, erro-padrão da média, desvio-padrão, nota mínima, mediana e nota máxima para cada um de seus componentes.
Tais estatísticas contemplam o total de estudantes concluintes da área de Zootecnia inscritos e presentes na prova do Enade 2023, tendo em vista agregações, ou por grandes regiões e o país como um todo, ou por categoria administrativa, organização acadêmica da IES e modalidade de oferta dos cursos.
Em relação aos gráficos de distribuição de notas, o intervalo considerado foi de 10 unidades, aberto à esquerda e fechado à direita, com exceção do primeiro intervalo, [ 0 ; 10 ], fechado em ambos os extremos.
Para os gráficos de distribuição das notas das questões discursivas, foram consideradas mais duas categorias : questão em branco⁶ e nota zero. ⁵ Uma definição dessas estatísticas pode ser encontrada no Glossário inserido no final desse relatório. ⁶ Nesse grupo estão incluídas também as respostas classificadas como nulas ou desconsideradas. 5.1 ESTATÍSTICAS BÁSICAS DA PROVA Esta seção apresenta estatísticas selecionadas e histogramas da nota geral ( 5.1.1 ) e de cada componente : Formação Geral ( 5.1.2 ) e Conhecimento Específico ( 5.1.3 ).
São também apresentadas estatísticas selecionadas de subpopulações, caracterizadas por grande região, categoria administrativa, organização acadêmica e modalidade de oferta. 5.1.1 ESTATÍSTICAS BÁSICAS GERAIS Na Tabela 5.1, são apresentadas as estatísticas básicas da prova, por grande região, dos estudantes concluintes de Zootecnia.
Como se pode verificar, a partir dos dados mostrados na Tabela 5.1, a média das notas da prova foi de 51,8 ( nas seções seguintes serão analisados os componentes de Formação Geral e de Conhecimento Específico ).
Os estudantes da região Norte obtiveram a média mais baixa ( 48,2 ), e os da região Sudeste e Sul alcançaram a média mais alta ( 54,1 ).
As demais médias foram 48,3, na região Centro-Oeste ; 51,7, na região Nordeste.
O desvio-padrão para o Brasil foi de 13,3, sendo o maior encontrado na região Sudeste ( 13,4 ), e o menor, na região Norte ( 12,0 ).
A região que obteve a maior nota máxima foi a região Sudeste ( 92,7 ), ao passo que a região que atingiu a menor nota máxima foi o Norte ( 80,9 ).
A mediana do Brasil, como um todo, foi igual a 52,0, sendo a maior mediana obtida na região Sudeste ( 54,8 ), e a menor, na região Centro-Oeste ( 47,6 ).
Considerando -se as notas segundo grande região, observa -se que existe diferença estatisticamente significativa ao nível de 95 \% entre as médias obtidas em todas as regiões, exceto entre Norte e Centro-Oeste ; Sul e Sudeste.
O comportamento das notas dos estudantes de todo o Brasil pode ser observado no Gráfico 5.1.
O Gráfico 5.1 apresenta um histograma com a distribuição das notas dos estudantes de Zootecnia de todo o Brasil : uma distribuição unimodal com a moda no intervalo ( 50 ; 60 ].
Na Tabela 5.2, são apresentadas informações referentes à média da nota final, desagregadas por categoria administrativa, organização acadêmica e modalidade de oferta dos cursos.
Os dados da Tabela 5.2 mostram que, em relação à categoria administrativa, os estudantes das IES públicas obtiveram média mais alta ( 52,1 ) que a média das IES privadas ( 47,8 ).
O desvio-padrão para as IES públicas foi de 13,2, indicando uma dispersão baixa das notas nesta categoria administrativa.
Já o desvio-padrão da IES privadas foi 13,3.
No que se refere à organização acadêmica, os estudantes de universidades obtiveram a média de 53,4, enquanto os de centros universitários, faculdades e CEFET/IF obtiveram médias de 46,7, 46,4 e 48,0, respectivamente.
Constata -se que existe diferença estatisticamente significativa, ao nível de 95 \%, entre as médias das organizações acadêmicas, exceto entre centro universitário e CEFET/IF ; faculdade e CEFET/IF ; faculdade e centro universitário.
A média da modalidade a distância ( 44,3 ) foi inferior à média da modalidade presencial ( 51,9 ), sendo a diferença entre as duas estatisticamente significativa, ao nível de 95 \%. 5.1.2 ESTATÍSTICAS BÁSICAS NO COMPONENTE DE FORMAÇÃO GERAL Na Tabela 5.3, são apresentadas as estatísticas básicas em relação ao componente da prova que avalia a Formação Geral dos estudantes concluintes.
A nota de Formação Geral é considerando -se os acertos nas nove questões objetivas e a nota média obtida na questão discursiva, sendo que a nota dessa questão obtida leva em conta 80 \% da nota relativa ao conteúdo e 20 \% da nota relativa ao desempenho linguístico.
Conforme os dados da Tabela 5.3, os estudantes de todo o Brasil obtiveram, no componente de Formação Geral da prova, desempenho médio de 47,7.
Quanto à variabilidade, o desvio-padrão das notas dos estudantes do Brasil foi 15,7.
A maior média foi obtida na região Sudeste ( 49,6 ), e a menor, na região Centro-Oeste ( 43,3 ).
As demais médias foram : 48,1, na região Nordeste ; 46,1, na região Norte ; 49,5, na região Sul.
Já o maior desvio-padrão foi obtido na região Nordeste e Sudeste ( 15,9 ), e o menor, na região Norte ( 14,8 ).
Os demais desvios foram : 15,4, na região Centro-Oeste ; e 15,1, na região Sul.
A maior nota no componente de Formação Geral da prova do Enade 2023 foi 90,2, obtida por pelo menos um estudante na região Nordeste.
A menor nota máxima foi obtida na região Centro-Oeste ( 78,3 ).
A mediana do Brasil, como um todo, foi de 48,2 ; sendo a menor mediana encontrada na região Centro-Oeste ( 43,4 ), e a maior, na região Sudeste e Sul ( 50,2 ).
A nota mínima nessa parte foi zero nas regiões Nordeste e Sudeste.
Considerando -se as notas, segundo grande região, observa -se que existe diferença estatisticamente significativa entre as médias das notas do componente de Formação Geral obtidas nas regiões, exceto entre as regiões Norte e Centro-Oeste ; Norte e Nordeste ; Sudeste e Nordeste ; Sul e Nordeste ; Sul e Sudeste.
No Gráfico 5.2, é possível verificar a avaliação do desempenho dos estudantes no componente de Formação Geral, a partir do histograma da distribuição das notas correspondentes.
Nesse gráfico, a distribuição é unimodal, com moda em ( 40 ; 50 ], mas com pouca diferença para o intervalo que o precede ( 50:60 ].
Na Tabela 5.4, são apresentadas as informações referentes ao desempenho dos concluintes no componente de Formação Geral, em diferentes agregações : categoria administrativa, organização acadêmica e modalidade de oferta.
A partir dos dados da Tabela 5.4, observa -se que há diferença estatisticamente significativa entre as médias dos tipos de categoria administrativa.
A maior média foi obtida por estudantes de IES públicas ( 48,1 ), e a menor, pelos de IES privadas ( 42,8 ).
Em relação às médias por tipos de organização acadêmica, encontra -se : 49,3 para as universidades ; 42,5 para os centros universitários, 40,3 para as faculdades e 44,5 para os CEFET/IF.
A nota máxima obtida pelas IES privadas ( 78,3 ) foi menor em comparação às IES públicas.
O desvio-padrão menor foi percebido nas IES privadas ; e, entre as organizações acadêmicas, esse valor foi menor para os centros universitários.
Já no que se refere à modalidade de oferta, os cursos presenciais apresentaram média de 47,8, enquanto a média dos cursos a distância foi de 41,5.
A diferença é, portanto, estatisticamente significativa. 5.1.3 ESTATÍSTICAS BÁSICAS DO COMPONENTE DE CONHECIMENTO ESPECÍFICO Na Tabela 5.5, são apresentadas as estatísticas básicas referentes ao componente de Conhecimento Específico da área de Zootecnia.
A nota do componente de Conhecimento Específico leva em conta as notas da parte objetiva e a média das notas obtidas na questão da parte discursiva da área.
De acordo com a Tabela 5.5, a média do desempenho dos estudantes do Brasil, para o componente de Conhecimento Específico da prova, foi de 53,1.
A maior média foi obtida na região Sul ( 55,7 ), e a menor, na região Norte ( 48,9 ).
As demais médias foram : 50,0, na região Centro-Oeste ; 52,9, na região Nordeste ; 55,6, na região Sudeste.
Quanto à variabilidade das notas, o desvio-padrão do Brasil foi 15,3 ; sendo o maior desvio-padrão observado na região Sudeste ( 15,6 ), e o menor, na região Norte ( 13,8 ).
Os demais desvios foram : 14,9, na região Centro-Oeste ; 15,0, na região Nordeste ; 15,2, na região Sul.
A mediana das notas dos estudantes de todo o Brasil foi de 52,8.
A maior mediana ocorreu na região Sudeste ( 56,3 ), e a menor, nas regiões Centro-Oeste e Norte ( 48,7 ).
As demais medianas foram : 53,3, na região Nordeste ; 55,1, na região Sul.
A nota máxima do Brasil, como um todo, foi 100,0 ; sendo obtida por, pelo menos, um estudante na região Sudeste.
As demais notas máximas foram : 93,7, na região Centro-Oeste ; 93,7, na região Nordeste ; 90,6, na região Norte ; 94,2, na região Sul.
A nota mínima nessa parte foi zero nas regiões Nordeste e Sudeste.
Observa -se que existe diferença estatisticamente significativa entre todas as regiões na análise das médias das notas do componente de Conhecimento Específico, exceto entre as regiões Norte e Centro-Oeste ; Sul e Sudeste.
Assim como nos Gráficos 5.1 e 5.2, mostrados anteriormente, no Gráfico 5.3, apresentado acima, foi possível ser realizada uma avaliação do desempenho de concluintes de Zootecnia, em relação ao componente de Conhecimento Específico, com um histograma da distribuição das notas correspondentes.
Esta também é uma distribuição unimodal, e o grupo modal é o ( 40 ; 50 ].
Na Tabela 5.6, é apresentada uma comparação dos resultados em relação à categoria administrativa, à organização acadêmica e à modalidade de oferta, levando -se em conta o desempenho de estudantes do componente de Conhecimento Específico da prova.
De acordo com os dados expostos na Tabela 5.6, no que se refere à organização acadêmica, a maior média foi das universidades ( 54,8 ), seguida dos CEFET/IF ( 49,2 ), e das faculdades ( 48,4 ) e, por fim, dos centros universitários ( 48,0 ).
O maior desvio-padrão foi o das faculdades ( 15,9 ).
As universidades obtiveram a maior mediana ( 54,3 ).
Já as notas máximas foram : 100,0 nas universidades, 90,6 nos centros universitários, 87,6 nas faculdades e 95,2 nos CEFET/IF.
A nota mínima foi zero para apenas nas universidades.
Observa -se que existe diferença estatisticamente significativa, ao nível de 95 \%, no componente de Conhecimento Específico entre as notas dos tipos de organização acadêmica, exceto entre centro universitário e CEFET/IF ; faculdade e CEFET/IF ; faculdade e centro universitário Quanto à categoria administrativa, existe diferença estatisticamente significativa entre as médias das IES públicas ( 53,4 ) e as das IES privadas ( 49,5 ).
Neste caso, assim como no componente de Formação Geral, a maior média foi obtida por estudantes de IES públicas de ensino.
Os cursos de educação presencial apresentaram maior média ( 53,2 ) em relação aos cursos a distância ( 45,2 ), sendo a diferença estatisticamente significativa. 5.2 ANÁLISE DAS QUESTÕES OBJETIVAS Esta seção apresenta estatísticas selecionadas e histogramas das questões objetivas dos componentes de Formação Geral ( 5.2.1 ) e de Conhecimento Específico ( 5.2.2 ).
São, também, apresentadas e comparadas as médias das subpopulações caracterizadas por grande região. 5.2.1 COMPONENTE DE FORMAÇÃO GERAL Na Tabela 5.7, são apresentadas as estatísticas básicas relativas às nove questões objetivas do componente da prova que abrange a Formação Geral dos estudantes.
De acordo com os dados da Tabela 5.7, em relação ao componente de Formação Geral, a média do Brasil foi 50,1.
A menor média foi encontrada na região Centro-Oeste ( 45,7 ), e a maior, na região Sudeste ( 52,7 ).
As demais médias foram : 50,3, na região Nordeste ; 47,3, na região Norte ; 51,9, na região Sul.
O desvio-padrão do Brasil foi 18,1 ; sendo o maior desvio-padrão encontrado na região Nordeste ( 18,6 ), e o menor desvio-padrão encontrado na região Norte ( 17,5 ).
A mediana do Brasil, como um todo, foi 55,6 e das regiões Sul, Nordeste, Centro-Oeste, Norte e Sudeste foram iguais a 55,6, 55,6, 44,4, 44,4 e 55,6, respectivamente.
Quanto à nota máxima, a região Sul alcançou 100,0 ; a Nordeste, 100,0 ; a Centro-Oeste, 88,9 ; a Norte, 100,0 ; e a Sudeste, 100,0.
Já as notas mínimas foi 0,0 em todas as regiões.
Na Tabela 5.8, exposta mais adiante, são apresentados o índice de facilidade e o índice de discriminação ( ponto-bisserial ) para cada uma das questões objetivas do componente de Formação Geral.
De acordo com os dados da Tabela 5.8, as questões objetivas do componente de Formação Geral, segundo o índice de facilidade, foram assim avaliadas : nenhuma teve o índice de facilidade classificado como “ Muito difícil ” ; nenhuma questão foi classificada como “ Muito fácil ” ; duas questões foram classificadas como “ Fácil ” ; cinco questões foram classificadas como"Médio ” ; e duas questões foram classificadas na categoria “ Difícil ”.
O índice de facilidade variou de 0,23 a 0,77.
Já em relação ao índice de discriminação, quatro questões foram classificadas com o índice “ Muito bom ” ; quatro com o índice “ Bom ”, uma questão com o índice “ Médio ” ; e nenhuma com o índice"Fraco ”.
O índice de discriminação variou de 0,27 a 0,46.
Na Tabela 5.9, é apresentada a distribuição das questões, levando -se em conta, simultaneamente, a classificação dos dois índices.
Os dados das tabelas 5.8 e 5.9, atestam que quatro questões que tiveram índice de discriminação “ Muito bom ” figuraram em dois níveis de dificuldade – “ Médio e Fácil ” : uma classificada na categoria “ Fácil ” ( Questão 6 ) ; três classificadas na categoria “ Médio ” ( questões 3, 4 e 8 ).
Em particular, a Questão 3 foi a que apresentou o maior poder discriminatório, com índice 0,46, considerado “ Médio ” em termos de facilidade, com uma proporção de 0,57 acertos, não tendo essa questão alcançado o máximo de acertos.
No Gráfico 5.4, para exemplificar, analisa -se o comportamento da Questão Objetiva 3 de Formação Geral.
Trata -se de uma questão considerada com nível “ Médio ” de dificuldade, em relação à facilidade, e que obteve o índice de discriminação 0,46 ( ” Muito bom ” ).
No Gráfico 5.4, cada uma das cinco curvas representa o percentual de respostas em determinada alternativa da questão, em função do número de acertos dos estudantes nessa parte da prova ( Formação Geral/múltipla escolha ), antes de possíveis eliminações pelo critério ponto-bisserial.
A curva em mostarda corresponde à alternativa B, a alternativa correta para essa questão.
Observa -se que, entre os estudantes com menor número de acertos, nessa parte do exame, a situação mais frequente foi a escolha da alternativa A ( em vermelho ), incorreta.
Por exemplo, entre os estudantes que acertaram até duas respostas, 31,3 \% escolheram a alternativa A ( em vermelho ) ; 16,2 \% escolheram a alternativa B ( em mostarda ) ; 14,0 \% escolheram a alternativa C ( em verde ) ; 26,0 \% escolheram a alternativa D ( em azul, gabarito ) ; e 11,5 \% escolheram a alternativa E ( em lilás ).
Entre os que acertaram duas respostas entre as questões de múltipla escolha de Formação Geral, 0,6 \% deixou a questão em branco e 0,6 \% marcou mais de uma alternativa, invalidando a questão.
À medida que o número de acertos aumenta, indicando desempenho melhor nessa parte da prova, aumenta, concomitantemente, a proporção de estudantes que selecionaram a alternativa correta B, atingindo 100 \% para os estudantes com oito acertos.
Essa análise permite verificar como a questão discriminou os grupos de desempenho, justificando -se o alto índice ( 0,46 ) obtido na questão.
Os gráficos relativos às demais questões de Formação Geral constam do Anexo IV. 5.2.2 COMPONENTE DE CONHECIMENTO ESPECÍFICO Na Tabela 5.10, são apresentadas as estatísticas básicas em relação às questões objetivas do componente de Conhecimento Específico da prova de Zootecnia por grande região.
Em relação ao componente de Conhecimento Específico da prova de Zootecnia, como será discutido ainda nesta seção, nove questões não foram usadas no cômputo das notas por terem sido descartadas pelo critério do ponto-bisserial e uma anulada.
Assim, as notas foram calculadas com base em 19 das 29 questões objetivas do componente de Conhecimento Específico.
Como se verifica a partir dos dados da Tabela 5.10, a média do Brasil para o componente de Conhecimento Específico da prova foi 52,1.
A menor média foi observada na região Norte ( 47,9 ), e a maior, na região Sudeste ( 54,8 ).
O desvio-padrão de todo o Brasil foi 15,9, sendo o maior desvio-padrão encontrado na região Sudeste ( 16,4 ), e o menor, na região Norte ( 14,5 ).
A mediana de todo o Brasil foi 52,6, o mesmo valor da mediana encontrada nas regiões Nordeste, Sudeste e Sul e superior às medianas das regiões Centro-Oeste e Norte ( 47,4 ).
A nota máxima da prova foi 100,0, obtida, nas questões objetivas do componente de Conhecimento Específico, por, pelo menos, um estudante da região Sudeste e, nas demais regiões, a nota máxima variou entre 89,5 ( Norte ) e 94,7 ( Centro-Oeste e Sul ).
A nota mínima de maior valor foi 10,5, obtida por, pelo menos, um estudante na região Norte.
Na Tabela 5.11, são apresentados os índices de facilidade e de discriminação ( ponto-bisserial ) das questões objetivas do componente de Conhecimento Específico, para os estudantes de Zootecnia.
A partir do índice de facilidade obtido, apresentado na Tabela 5.11, pode -se concluir que a maioria das questões objetivas da prova foram classificadas no índice “ Difícil ” ( 11 ).
Sete questões foram classificadas no índice"Fácil ” ; nove, no índice “ Médio ” ; nenhuma, no índice “ Muito fácil ” ; e uma questão foi alocada no índice “ Muito difícil ”.
Quanto aos índices de discriminação das questões objetivas do componente de Conhecimento Específico da prova, obtém -se como resultado a seguinte classificação : nenhuma questão foi alocada no índice “ Muito bom ” e Onze, no “ Bom ”.
Dessa forma, 11 das 28 questões foram classificadas com os índices de discriminação “ Bom ” e “ Muito bom ”.
Entre as demais, oito delas foram alocadas no índice “ Médio ” e Na Tabela 5.12, é apresentada a distribuição das questões considerando -se, simultaneamente, a classificação dos dois índices.
Os dados das tabelas 5.11 e 5.12 atestam que 11 questões que tiveram índice de discriminação “ Bom ” figuraram em três níveis de dificuldade — “ Difícil","Fácil"e"Médio ” : quatro classificadas na categoria “ Fácil ” ; quatro, na categoria “ Médio ” ; três, na categoria “ Difícil ” ; nenhuma, na categoria “ Muito fácil ” e nenhum, na categoria “ Muito difícil ”, As questões que tiveram o índice de discriminação “ Fraco ” configuraram -se com índices de facilidade “ Muito fácil ” ( nenhuma questão ), “ Fácil ” ( nenhuma questão ), “ Médio ” ( duas questão ), “ Difícil ” ( seis questões ) e “ Muito Difícil ” ( uma questões ).
Entre as questões que alcançaram os maiores índices de discriminação, “ Muito bom ”, uma delas, a questão 36 e 38, teve o maior índice de discriminação ( 0,39 ).
Tal questão foi classificada como “ Fácil ”, com índice de facilidade igual a 0,61 e 0,62, respectivamente.
A questão 22 foi a mais “ Difícil ” entre as 28 questões específicas válidas, com baixo índice de facilidade, apenas 0,14 de proporção de acertos.
Essa questão apresentou poder discriminatório fraco, 0,13, o que comprova ter sido a mais “ Difícil ” para os estudantes.
Destacam -se, também, com baixo índice de facilidade, as questões 19 e 22, cujos índices de facilidade, em termos proporcionais, correspondem, respectivamente, a 0,19 e 0,14 de estudantes que responderam acertadamente.
Essas duas questões também tiveram poder discriminatório “ Fraco"e"Fraco ”, com índices 0,08 e 0,13, respectivamente.
Destaca -se que as questões 16, 17, 19, 22, 24, 25, 26, 31 e 32 foram consideradas inadequadas pelo critério ponto-bisserial.
Por isso, foram eliminadas do cômputo da nota final.
A título de exemplo das análises do comportamento das questões objetivas, no Gráfico 5.5, analisa -se a questão 23 do componente de Conhecimento Específico.
No Gráfico 5.5, cada uma das cinco curvas representam o percentual de respostas em determinada alternativa da questão, em função do número de acertos dos estudantes nessa parte da prova ( componente de Conhecimento Específico ), antes de possíveis eliminações pelo critério ponto-bisserial.
A alternativa correta B, representada no gráfico pela curva em mostarda, foi escolhida em maiores proporções pelos estudantes com desempenho melhor nessa parte da prova.
Já as alternativas incorretas, também denominadas distratores, foram selecionadas, principalmente, por aqueles com notas mais baixas.
A Questão 23 foi considerada, pelas respostas dos estudantes avaliados na prova, como uma questão com índice de facilidade “ Médio ” ( 0,56 ), ou seja, 56,0 \% dos estudantes assinalaram, acertadamente, a opção B correspondente ao gabarito.
Seu índice de discriminação foi igual a 0,35, classificado como “ Bom ”, o maior valor de discriminação.
Observa -se que a soma não é 100 \%, em razão das questões não respondidas ou com mais de uma opção marcada.
Cerca de 100,0 \% daqueles com nota zero marcaram a opção em branco ou dupla marcação.
A proporção de estudantes que selecionaram a resposta correta B aumenta gradativamente, chegando a atingir 100 \% para 22 acertos, enquanto a proporção dos que escolheram alternativas incorretas decai, a partir de sete acertos, em função do número de acertos nessa parte da prova.
Os gráficos relativos às demais questões do componente de Conhecimento Específico constam no Anexo IV.
Relatório Síntese de Área – ENADE/2023 5.3 ANÁLISE DAS QUESTÕES DISCURSIVAS Esta seção apresenta estatísticas selecionadas e histogramas das questões discursivas de Formação Geral ( 5.3.1 ) e de Conhecimento Específico ( 5.3.2 ).
Expõe e compara, também, as médias de subpopulações, caracterizadas por grande região.
Cumpre notar que nem todas as questões passam pelo crivo dos corretores.
Após a digitalização, as respostas passam por um pré-processamento de identificação automático.
Algumas dessas questões são definidas como um “ branco automático ” ( que não são enviadas para os professores para correção, a não ser que, na segunda fase, a de inspeção visual, seja constatada uma resposta ).
O critério usado nesse pré-processamento consiste em verificar, para cada questão, a partir da quantidade de pixels, o que, provavelmente, é uma resposta em branco.
A partir da separação, os documentos que foram considerados em branco são enviados para uma verificação visual, feita por uma equipe de colaboradores devidamente capacitados.
Desse modo, qualquer erro, presente no pré-processamento, é encontrado, e as questões preenchidas são enviadas para a correção dos avaliadores.
Eventualmente, algumas respostas em branco não são detectadas nesse pré-processamento automático.
A presença de pequenos riscos, ou manchas, pode fazer o programa não identificar uma questão como um possível branco.
Essas são enviadas para a correção, mas são classificadas como “ branco ” pelos avaliadores, no momento da correção. 5.3.1 COMPONENTE DE FORMAÇÃO GERAL As análises dos resultados de desempenho dos estudantes de Zootecnia, na questão discursiva relativa à Formação Geral, encontram -se na Tabela 5.13 e no Gráfico 5.6, cujos dados são analisados de forma \textbf{sequenciada}.
No caso da Tabela 5.13, analisa -se a média das notas da questão discursiva de Formação Geral, que leva em conta as correções realizadas por duas bancas : uma que analisa e pontua aspectos relativos ao conteúdo das respostas, seguindo o padrão de respostas divulgado pelo Inep ( ver Anexo VIII ) ; e outra que avalia o desempenho linguístico expresso pela redação das respostas, seguindo um padrão de respostas específico da área de Língua Portuguesa, detalhado na seção 5.3.1.5 deste relatório.
No cálculo da nota da questão discursiva, consideram -se 80 \% relativos à nota da correção de conteúdo e 20 \% referentes à nota de desempenho linguístico.
Na Tabela 5.13, observa -se que os estudantes de todo o Brasil obtiveram, em Formação Geral, média de 43,3 nas questões discursivas.
Pode -se notar, também, que o desvio-padrão nesse conjunto de questões foi de 23,9.
A maior média foi obtida na região Sul ( 45,0 ), e a menor, na região Centro-Oeste ( 38,8 ).
A mediana de todo o Brasil, neste componente, foi 43,7 e, na região Sul ( 46,5 ), foi encontrada a maior mediana.
A nota máxima ( 100,0 ) foi obtida por, pelo menos, um estudante da região Norte, Sudeste e Sul.
A nota mínima ( 0,0 ) foi a mesma em todas as regiões do Brasil.
No Gráfico 5.6, está representada a distribuição das notas na questão discursiva do componente de Formação Geral.
Como mostra o Gráfico 5.6, a moda da distribuição ocorre no intervalo ( 50 ; 60 ], com frequência de 15,0 \% dos participantes, sendo que, no intervalo [ 0 ; 10 ], inclui -se, além da nota zero, a frequência de estudantes que deixaram esse tipo de questão em branco.
Na sequência, os resultados verificados para a questão discursiva de Formação Geral serão apresentados, estabelecendo -se relações com os conteúdos abordados em cada uma delas.
Os comentários das bancas de docentes avaliadores, a respeito do observado na correção das respostas dos estudantes, suas impressões e conclusões serão apresentadas junto à análise de cada questão.
Cumpre esclarecer que, tendo em vista que a questão discursiva de Formação Geral foi padronizada, ou seja, foi a mesma em todas as provas, os comentários da banca são os mesmos para todas as carreiras acadêmicas, sendo direcionados a todos os estudantes que participaram do Enade 2023.
A seguir, serão analisados o desempenho linguístico dos estudantes da área de Zootecnia na questão discursiva de Formação Geral do Enade 2023, apresentando -se os resultados obtidos e os comentários da banca de correção sobre a questão. 5.3.1.1 ANÁLISE DE CONTEÚDO DA QUESTÃO DISCURSIVA DO COMPONENTE DE FORMAÇÃO GERAL Os dados de Zootecnia, obtidos a partir das respostas à questão discursiva do componente de Formação Geral, encontram -se na Tabela 5.14 e no Gráfico 5.7, exposto após a apresentação dos dados da referida tabela.
Verifica -se, a partir dos dados da Tabela 5.14, que, no conteúdo da questão discursiva do componente de Formação Geral, os estudantes de todo o Brasil obtiveram média de 38,6.
A maior média foi obtida na região Sul ( 40,4 ), e a menor, na região Centro-Oeste ( 34,1 ).
Quanto à variabilidade das notas, o desvio-padrão de todo o Brasil foi 27,5.
O menor desvio-padrão foi obtido na região Norte e Sul ( 27,0 ), e o maior desvio-padrão foi obtido pela região Centro-Oeste ( 27,7 ).
A mediana do Brasil, como um todo, foi 37,5, sendo superior, igual e igual às das regiões Centro-Oeste ( 32,0 ), Norte ( 37,5 ) e Sudeste ( 37,5 ) e inferior à mediana das regiões Nordeste e Sul ( 40,0 ).
A nota No Gráfico 5.7, apresenta -se a distribuição das notas de conteúdo da questão discursiva do componente de Formação Geral.
O percentual de participantes que deixaram a prova em branco ou zeraram a questão discursiva na avaliação de conteúdo foi 16,4 \%.
O Gráfico 5.7 mostra que a moda dessa distribuição ocorre no intervalo [ 0 ; 10 ], com uma frequência de 20,6 \% dos participantes. 5.3.1.2 COMENTÁRIOS SOBRE A CORREÇÃO DE CONTEÚDO DAS RESPOSTAS À QUESTÃO DISCURSIVA A questão discursiva do componente de Formação Geral foi composta por um texto motivador e por dois itens ( a e b ).
O texto motivador apresentou um comentário sobre uma publicação do Instituto Brasileiro de Geografia e Estatística ( IBGE ) que tratava dos indicadores sobre a realidade social brasileira, apontando o crescimento do nível de desigualdade socioeconômica em 2021 e concluindo que a pobreza teve aumento recorde nesse período.
A questão foi dividida em dois itens.
O primeiro item solicitava que o respondente estabelecesse uma relação entre o perfil da população mais afetada pelas desigualdades sociais e as situações de risco socioambiental.
Já o segundo item pedia a proposição de duas possíveis ações a serem desenvolvidas em bairros periféricos a fim de minimizar os impactos socioambientais.
A questão foi pertinente, apresentou linguagem acessível e abordou conteúdo basilar para Formação Geral ; no entanto, apresentou um nível de dificuldade alto, dadas as competências intrincadas que demandou.
Além disso, o tema foi atual e de grande relevância para a sociedade brasileira, visto que abordou assunto de impacto social e de provável domínio dos concluintes de curso superior.
O texto motivador foi adequado e coerente com o que se pediu na questão, tratou do aumento da desigualdade social e da descrição do perfil da população com base em critérios socioeconômicos.
Com relação às habilidades cognitivas, observou -se que as seguintes foram necessárias para responder à questão proposta : capacidade de expressar ideias de maneira clara, coesa e argumentativa por meio da produção de textos ; capacidade de inter-relacionar objetos de conhecimento nas diferentes áreas ; e capacidade de criar estratégias de ação, selecionar métodos e formular propostas.
Entre os saberes necessários para o desenvolvimento das respostas, os respondentes deveriam conhecer o meio social e/ou geográfico da população brasileira de perfil socioeconômico pobre ou extremamente pobre e deveriam saber relacionar esse meio com situações de risco socioambiental que afligem principalmente esse perfil populacional.
Portanto, deveriam, também, conhecer uma variedade de riscos socioambientais a que esse perfil populacional está sujeito, por conta do meio social e/ou geográfico em que vive.
A questão, em seu item b ), por sua vez, demandou o domínio de um rol de propostas minimamente factíveis para reduzir tais situações de risco, envolvendo distintos atores sociais.
Em relação ao padrão de resposta, esperava -se para cada um dos itens avaliados que : a ) O estudante explicasse a relação entre desigualdade social e riscos socioambientais percorrendo o seguinte trajeto teórico-argumentativo : • descrever o(s) risco(s) socioambiental(is) que afeta(m) a população com perfil pobre ou extremamente pobre, tais como inundações, deslizamentos de terra, contaminação ambiental, maior vulnerabilidade em relação a doenças, discriminação social e criminalidade, entre outros ; • descrever o meio ( geográfico ou social ), como morros, áreas próximas a rios, mangues, espaços marcados pela violência urbana, entre outros, nas cidades brasileiras, como fator de criação ou potencialização das condições de vulnerabilidade e risco aos pobres ou extremamente pobres ; • estabelecer nexo causal entre a pobreza ou extrema pobreza, o meio ( geográfico ou social ) urbano e a situação de risco socioambiental, relacionando esses aspectos ao perfil socioeconômico da população.
Portanto, esperava -se que o estudante apresentasse consistência argumentativa quanto ao que foi solicitado, notadamente quanto aos riscos socioambientais, à descrição do meio e ao estabelecimento do nexo causal entre pobreza, meio e risco socioambiental. b ) O estudante apresentasse propostas pertinentes, factíveis e bem desenvolvidas que envolvessem ação governamental ( federal, estadual/distrital ou municipal ) e participação da comunidade a fim de minimizar riscos socioambientais, como, por exemplo : • implantar obras de infraestrutura urbana que envolvessem contenção de morros e encostas ; saneamento básico : serviços regulares de abastecimento de água, coleta e tratamento de esgotos, coleta e manejo de resíduos sólidos, manejo de águas pluviais, entre outras ; • desenvolver políticas sociais, como segurança pública, provisão de moradia adequada, com custo acessível, e regularizada do ponto de vista fundiário ; educação de qualidade ; segurança pública ; atendimento psicossocial, com especial atenção para grupos em situação de vulnerabilidade ; entre outras ; • garantir o acesso dessas comunidades a equipamentos sociais e culturais, bem como a espaços públicos inclusivos e a áreas verdes ; • elaborar planos estratégicos participativos urbanos e ambientais a serem desenvolvidos nas comunidades em situação de vulnerabilidade.
No tocante ao subitem b ), esperava -se que o respondente apresentasse coerência propositiva, englobando os diferentes atores, quais sejam : o governo, em suas diversas instâncias, e a comunidade, também em suas diversas organizações, almejando conter os riscos socioambientais. 5.3.1.5 CONSIDERAÇÕES FINAIS SOBRE A CORREÇÃO DE CONTEÚDO DA QUESTÃO DISCURSIVA DO COMPONENTE DE FORMAÇÃO GERAL De uma maneira geral, os respondentes apresentaram um nível aquém do esperado em termos de formação geral, destacando -se, neste âmbito, os aspectos a seguir.
Nas respostas relacionadas ao item a ), percebeu -se que, quanto aos riscos, a maioria dos respondentes não se limitou a descrevê -los, ou mesmo a estabelecer um rol deles.
No caso daqueles que descreveram as diversas situações de risco, em sua maioria, as respostas se enquadraram naquilo que estava previsto no padrão de resposta, apresentando os riscos derivados da insegurança relativa ao meio geográfico em que habitam, tais como alagamentos, deslizamentos e doenças por contato com um meio insalubre.
Quanto à descrição do meio, notou -se que houve uma maior aderência por parte dos respondentes, pois a maior parte das respostas, excluindo -se as que foram deixadas em branco, descreveu o meio de alguma forma.
Devido à grande variedade de possíveis descrições do meio geográfico e/ou social em que vive a população brasileira de perfil pobre ou extremamente pobre, o padrão de respostas deixava aberta a possibilidade de considerar respostas variadas em relação às citadas como exemplo em sua redação.
E, de fato, entre as respostas encontradas, verificaram -se muitas que foram consideradas certas, por estarem dentro do perfil previsto como correto, mas que eram distintas dos exemplos específicos apresentados no padrão de respostas.
O mesmo se deu quanto ao item b ) : houve uma grande variedade de proposições para além dos campos temáticos específicos previstos no padrão de respostas, mas que foram consideradas corretas porque atendiam à demanda do item.
Para esse item, o padrão de respostas demandava que as proposições fossem pertinentes ao tema e minimamente desenvolvidas.
No entanto, a maior parte dos respondentes não apresentou proposições desenvolvidas e muitos não fizeram proposições pertinentes ao tema.
O item também demandava que as proposições abrangessem a ação governamental e a participação da comunidade, mas a maior parte dos respondentes envolveu em suas proposições apenas os governos em suas mais diversas instâncias, não abrangendo o papel participativo da comunidade.
Um pequeno grupo de respostas se limitou a apresentar só uma proposição.
Dessa forma, o desempenho dos estudantes foi, em geral, muito baixo.
O número de questões deixadas em branco foi considerável, bem como o quantitativo de respostas muito curtas, ocupando menos de cinco das 15 linhas oferecidas para a resposta.
Muitos responderam a apenas um dos dois itens.
A grande maioria dos que buscaram responder ao que a questão demandava, completamente ou parcialmente, não souberam se expressar de modo correto, objetivo e coerente.
Não poucos respondentes repetiram longos trechos do enunciado em suas respostas, demonstrando um baixo nível de formação geral e, também, baixa disposição para a realização da prova.
Muitos respondentes foram desatentos e escreveram a resposta à questão discursiva de componente de Conhecimento Específico no campo dedicado à questão discursiva de Formação Geral, tendo sido incluídos dentre os respondentes que fugiram ao tema.
Uma quantidade proporcionalmente pequena, mas grande numericamente, foi a de estudantes que apresentaram respostas que desrespeitavam os direitos humanos, em sua maioria propondo o controle reprodutivo sobre a população pobre ou extremamente pobre como solução governamental para as situações de risco socioambiental comumente vivenciadas pelas populações brasileiras pobres ou extremamente pobres.
Alguns estudantes também apontaram uma relação necessária entre pobreza e ignorância ou mesmo entre pobreza e violência.
Entre os estudantes que responderam com maior dedicação, buscando dar conta de toda a questão, houve um número significativo que aparentou ter grandes dificuldade de expressão textual, dando a entender que seu baixo domínio da linguagem verbal comprometia a capacidade de produzir respostas adequadas.
Muitas respostas não apresentaram o mínimo de lógica ou não buscaram estabelecer relações entre os conteúdos mobilizados.
Portanto, em sua maioria, os estudantes não apresentaram a capacidade de expressar ideias de maneira clara, coesa e argumentativa por meio da produção de textos, uma habilidade fundamental ao concluinte do Ensino Superior.
Um outro perfil de destaque foi o de estudantes que se dedicaram a pontuar elementos válidos para responder à questão sem utilizar os conectivos, as relações causais.
Também merecem destaques os respondentes que falaram sobre aspectos históricos da pobreza e a descrição dos tipos sociais mais envolvidos em situação de pobreza ou extrema pobreza, fugindo de buscar estabelecer uma relação causal com as situações de risco socioambiental, portanto, fugindo do que a questão solicitava.
Essas situações demonstraram um perfil majoritário com baixa capacidade de produzir para o item de Formação Geral uma resposta adequada esperada para estudantes concluintes do Ensino Superior.
Os estudantes com perfis de respostas minimamente adequadas ou boas corresponderam ao segundo perfil majoritário.
Em sua maior parte, esse perfil de estudantes respondeu bem ao subitem a ), mas não respondeu bem ou subitem b ), ou vice-versa.
O perfil minoritário foi o de estudantes de nível muito bom a ótimo.
Esses com competência para descrever objetivamente os pontos essenciais da questão, estabelecendo o nexo causal entre eles.
Quanto à relação com o conteúdo, a maior parte dos respondentes soube pensar nos elementos e nas imagens do problema levantado na questão, mas não soube fazer as ligações que estimulam a relação entre esses elementos.
Portanto, demonstraram uma baixa capacidade de inter-relacionar objetos de conhecimento nas diferentes áreas, outra habilidade esperada de um concluinte do Ensino Superior.
Uma grande parte dos estudantes vê o saneamento básico precário como a maior fonte de riscos socioambientais para a população pobre ou extremamete pobre.
As descrições do meio físico, como as situações de moradias em encostas de morros e beiras de rio, também constituíram entre as mais comuns.
Em menor número, mas ainda com relativa frequência, apareceram os casos de estudantes que apontaram o meio como lugar de domínio territorial de criminosos ou os casos que indicaram o meio social como opressivo aos negros.
Quanto às proposições, os estudantes, em sua maioria, não souberam criar estratégias de ação, selecionar métodos e formular propostas.
Uma quantidade considerável de estudantes apresentou proposições muito genéricas como “ construir escolas ”.
A maioria não apontou conexão entre as propostas e a minimização dos riscos socioambientais, como se a ação proposta já evidenciasse o que ela está solucionando, assim também indicando uma grande concentração de estudantes que não aparenta dominar habilidades importantes para concluintes do Ensino Superior. 5.3.1.6 ANÁLISE DE LÍNGUA PORTUGUESA DA QUESTÃO DISCURSIVA DO COMPONENTE DE FORMAÇÃO GERAL Os dados de Zootecnia obtidos a partir das respostas à questão discursiva do componente de Formação Geral, no que tange ao desempenho linguístico, são apresentados na Tabela 5.15.
Os dados da Tabela 5.15 mostram que, no quesito desempenho linguístico, os estudantes, de todo o Brasil, obtiveram média de 62,0.
A maior média em relação ao desempenho linguístico foi obtida na região Sudeste ( 63,5 ), e a menor, na região Centro-Oeste ( 57,6 ).
Quanto à variabilidade das notas, o desvio-padrão de todo o Brasil foi 21,4.
O menor desvio-padrão foi obtido na região Norte ( 19,8 ), e o maior desvio-padrão foi obtido na região Centro-Oeste ( 24,4 ).
A mediana das notas de Língua Portuguesa foi 67,5 para o Brasil, como um todo, a mesma obtida em todas as regiões, exceto no Norte ( 65,0 ).
A nota máxima para todo o Brasil foi 100,0, com, pelo menos, um estudante com essa nota em todas as regiões.
A nota mínima foi 0,0 em todas as regiões do país.
No Gráfico 5.8, mostra -se a distribuição das notas de desempenho linguístico do componente de Formação Geral.
O Gráfico 5.8 demonstra que 6,4 \% dos participantes correspondem ao intervalo [ 0 ; 10 ].
A moda dessa distribuição ocorre no intervalo ( 70 ; 80 ], com uma frequência de 24,9 \% dos participantes. 5.3.1.7 COMENTÁRIOS SOBRE A CORREÇÃO DE LÍNGUA PORTUGUESA DAS RESPOSTAS À QUESTÃO DISCURSIVA Os comentários feitos neste item a respeito da correção de Língua Portuguesa das respostas emitidas para a questão discursiva do componente de Formação Geral estão agrupados em três pontos : avaliação da questão, padrão de resposta utilizado na avaliação e análise das respostas dadas pelos estudantes.
A linguagem da questão e do texto motivador foi considerada condizente com o domínio linguístico que se espera de estudantes concluintes de cursos em nível de graduação, sendo -lhes, portanto, acessível.
Do ponto de vista linguístico e textual, o nível da questão foi considerado mediano.
Por meio de um comando claro e objetivo, a questão solicitava que o estudante explicasse determinada relação, exposta no texto motivador, e apresentasse duas propostas para atenuar o problema discutido.
Para compreender o que deveria ser feito, o estudante necessitava ter domínio da leitura e interpretação textual, de modo a entender o enunciado, estabelecer a relação proposta na questão e apontar possíveis ações governamentais e da comunidade para resolver os problemas elencados.
O fato de a proposta textual ser apresentada em dois itens ( a e b ) separadamente levou muitos estudantes a produzirem o texto em dois parágrafos bem definidos e, na maioria dos casos, sem conexão textual entre ambos, ou seja, respondendo à questão em dois ou mais parágrafos que indicavam a letra do item ao qual a resposta se dirigia.
Foram relativamente poucas as provas em que os estudantes respondiam à questão em um texto que articulasse textualmente os dois itens.
Essa formatação dificultou a avaliação do quesito Aspectos Textuais no que se refere à organização textual.
Como a questão demandava certa capacidade de análise, a resposta exigia, em termos textuais, que o estudante fosse além do texto motivador, o que, em parte, passou a identificar a qualidade textual da resposta, ainda que tomasse como base o próprio enunciado da questão e empregasse dados do texto motivador.
Textos com qualidade mediana, em geral, consistiram na retextualização da questão e/ou de enunciados do texto motivador, com alguma autonomia linguística para elaborar argumentação superficial.
A simples cópia — às vezes, com erros ortográficos ou morfossintáticos — da questão ou de enunciados do texto motivador, praticamente sem autonomia textual ou argumentativa, foi um recurso característico de respostas de baixa qualidade em termos linguísticos, porém não penalizado na avaliação do desempenho linguístico.
Em relação ao padrão de resposta utilizado na avaliação, foram considerados os seguintes aspectos : ( 1 ) aspectos ortográficos ; ( 2 ) aspectos textuais ; ( 3 ) aspectos morfossintáticos e de propriedade vocabular adequada à modalidade escrita formal da Língua Portuguesa, exigida pela situação comunicativa.
A seguir, são apresentados cada um dos aspectos avaliados : No que concerne à análise das respostas dadas pelos estudantes, especificamente, em relação ao desempenho linguístico, foram observados muitos deslizes gramaticais no que diz respeito à acentuação e grafia correta das palavras, embora, de maneira geral, os estudantes tenham demonstrado domínio da ortografia vigente.
Quanto aos aspectos textuais, verificou -se predominância de textos organizados sem a presença das três partes esperadas para um texto dissertativo-argumentativo : introdução ; desenvolvimento e conclusão.
A argumentação foi predominantemente rasa e pouco elaborada, havendo frequentes desvios de pontuação, nomeadamente quanto ao emprego da vírgula.
Por fim, quanto aos aspectos morfossintáticos, observou -se um bom domínio das regras de regência e concordância na maior parte dos textos.
A adequação lexical também foi pertinente a uma redação formal na maior parte das vezes.
Em relação ao item a ), destacaram -se como respostas a relação entre comunidades negras periféricas, além de ribeirinhos e indivíduos em vulnerabilidade socioeconômica como sendo o perfil da população brasileira mais propenso a riscos socioambientais.
No que diz respeito ao item b ), foram encontradas com maior frequência propostas relacionadas a atuações governamentais, familiares e sociais diversas, como a necessidade de maiores investimentos e ofertas de trabalho na periferia. 1 ) Aspectos ortográficos – referem -se ao domínio das convenções ortográficas da norma-padrão da Língua Portuguesa, como : grafia de vogais e consoantes, uso de maiúsculas e minúsculas, emprego do hífen e da acentuação gráfica, segmentação de palavras. 2 ) Aspectos textuais – referem -se ao domínio dos procedimentos de estruturação textual : organização interna dos períodos, emprego de conectores para a articulação lógica entre os períodos e entre os parágrafos, emprego de marcas de referenciação lexical e pronominal ; e utilização dos sinais de pontuação que contribuem para a organização lógica da frase. 3 ) Aspectos morfossintáticos e de propriedade vocabular – referem -se ao domínio das regras de caráter morfossintático estabelecidas como modelares do ponto de vista da modalidade escrita formal da norma-padrão da Língua Portuguesa : concordância nominal e verbal, regência nominal e verbal, colocação pronominal, flexão nominal e verbal, correlação entre tempos e modos verbais, ausência de marcas de oralidade.
Outra competência combinada a esta na análise das respostas foi a seleção vocabular adequada à modalidade escrita formal da Língua Portuguesa, tendo em vista a \textbf{interseção} entre as duas do ponto de vista das exigências do registro formal da modalidade escrita da norma-padrão.
No que concerne à análise das respostas dadas pelos estudantes, especificamente, em relação ao desempenho linguístico, foram observados muitos deslizes gramaticais no que diz respeito à acentuação e grafia correta das palavras, embora, de maneira geral, os estudantes tenham demonstrado domínio da ortografia vigente.
Quanto aos aspectos textuais, verificou -se predominância de textos organizados sem a presença das três partes esperadas para um texto dissertativo-argumentativo : introdução ; desenvolvimento e conclusão.
A argumentação foi predominantemente rasa e pouco elaborada, havendo frequentes desvios de pontuação, nomeadamente quanto ao emprego da vírgula.
Por fim, quanto aos aspectos morfossintáticos, observou -se um bom domínio das regras de regência e concordância na maior parte dos textos.
A adequação lexical também foi pertinente a uma redação formal na maior parte das vezes.
Em relação ao item a ), destacaram -se como respostas a relação entre comunidades negras periféricas, além de ribeirinhos e indivíduos em vulnerabilidade socioeconômica como sendo o perfil da população brasileira mais propenso a riscos socioambientais.
No que diz respeito ao item b ), foram encontradas com maior frequência propostas relacionadas a atuações governamentais, familiares e sociais diversas, como a necessidade de maiores investimentos e ofertas de trabalho na periferia.
Houve a ocorrência de respostas fracas e incompletas.
O item a ) foi respondido, na maior parte dos textos, de modo discursivo, enquanto o item b ) figurou, por vezes, como uma lista de ações que poderiam resolver o problema.
Em alguns casos, não tão raros, foram produzidos textos nos quais os estudantes se recusaram explicitamente a produzir uma resposta adequada, alegando que a questão era pertinente a uma área de formação diferente da sua.
As respostas encontradas e mesmo a ausência de respostas nos permitem tecer algumas reflexões.
Há constante falta de atenção na leitura do enunciado da questão, o que levou alguns estudantes a responderem à questão específica em local diverso do esperado, prejudicando a resposta.
Por vezes, alguns candidatos produziram as respostas para as duas questões ( Formação Geral e Conheciomento Específico ) no mesmo espaço, prejudicando o desempenho textual e argumentativo.
Seja por desatenção, seja por descuido, a dificuldade de interpretação é sintomática e deve ser considerada.
Desvios ortográficos, morfossintáticos e de pontuação frequentes apontam para lacunas no que concerne à alfabetização e ao desenvolvimento da escrita ainda nos anos do ensino fundamental.
Contudo, o número baixo de textos que tenham obtido notas muito baixas, como 1,0 e 2,0, revela que, entre os estudantes avaliados, não havia pessoas com níveis extremamente baixos de alfabetização e letramento.
Problemas quanto à organização do texto dissertativo-argumentativo, com ausência de introdução, desenvolvimento e conclusão, mesmo que minimamente, apontam para a necessidade de um trabalho mais profundo com a escrita em cursos superiores.
A argumentação rasa ou mesmo o desinteresse por temas sociais revelam o fato de que alguns dos cursos estudados não trabalham o senso crítico do estudante, não o conduzindo a uma vivência plenamente crítica e democrática em sociedade.
A dificuldade em responder ao item b ) da questão demonstrou uma fragilidade teórica na formação, principalmente no que concerne à atuação do profissional em relação a questões de cunho socioeconômico e ao atendimento ao público mais fragilizado socioeconomicamente.
A baixa frequência de desvios morfossintáticos pode ser justificada pela própria competência linguística do estudante, demonstrando que um trabalho mínimo com questões gramaticais nos cursos superiores poderia sanar grande parte do problema.
Relativamente poucos estudantes demonstraram incapacidade quanto à escolha de palavras adequadas para a produção de um texto formal, apresentando competência genérica, na maior parte das vezes, satisfatória.
Textos que obtiveram nota máxima foram raros, o que aponta para a necessidade de reforçar a importância do trabalho com a leitura e escrita em âmbito formal com os estudantes dos cursos avaliados. 5.3.1.8 CONSIDERAÇÕES FINAIS SOBRE A CORREÇÃO DE LÍNGUA PORTUGUESA DAS RESPOSTAS À QUESTÃO DISCURSIVA DO COMPONENTE DE FORMAÇÃO GERAL A análise global em relação ao desempenho linguístico dos cursos avaliados no Enade 2023 revelou um desempenho mediano dos estudantes avaliados.
Quanto aos aspectos textuais e argumentativos, verificou -se um nível baixo de conhecimento e mesmo de interesse, por parte de muitos estudantes, em refletir sobre os problemas sociais envolvendo pessoas em situação de vulnerabilidade e fenômenos de risco socioambiental.
A maioria dos estudantes não conseguiu apresentar uma análise crítica, tampouco responder conforme solicitado pela questão.
Alguns estudantes não foram atentos às informações apresentadas nos textos nem ao que os enunciados solicitavam.
Alguns demonstraram o domínio sobre o assunto ou, no mínimo, o comprometimento e o interesse para responder as questões.
Entende -se que o tempo, a quantidade de questões e a pressão podem influenciar no momento de elaborar a resposta.
Além disso, alguns estudantes de cursos específicos não se sentiram motivados a falar de um assunto julgado como “ social ” e não pertinente à sua área de atuação profissional, o que é um forte representativo de vieses ideológicos que perpassam determinadas formações.
Cabe ressaltar que a questão abordava um conteúdo atual e relevante a ser considerado por um indivíduo plenamente capaz de agir em sociedade de forma crítica, de qualquer área de formação em nível superior.
Destaca -se que estudantes com aparente domínio da área de saúde apresentaram inclinação crítica a abordar o tema de forma bastante humanizada.
Quanto aos aspectos ortográficos e gramaticais de natureza morfossintática, observam -se muitas lacunas trazidas desde a educação básica.
Entende -se que a universidade não pode sanar todos os problemas de cunho ortográfico e gramatical, mas é preciso salientar a importância de disciplinas voltadas à leitura e produção de texto serem incluídas como disciplinas obrigatórias nos cursos avaliados. 5.3.2 COMPONENTE DE CONHECIMENTO ESPECÍFICO As análises dos resultados de desempenho dos estudantes de Zootecnia, na questão discursiva relativa ao componente de Conhecimento Específico, encontram -se na Tabela 5.16 e no Gráfico 5.9. 5.3.2.1 ANÁLISE DA QUESTÃO DISCURSIVA DO COMPONENTE DE CONHECIMENTO ESPECÍFICO A Tabela 5.16 apresenta dados sobre a questão discursiva do componente de Conhecimento Específico.
A partir dos dados registrados na Tabela 15.16, observa -se que a média foi mais baixa se comparada à média obtida na questão discursiva do componente de Formação Geral.
Enquanto, no componente de Formação Geral, a média para estudantes de Zootecnia de todo o Brasil foi 43,3 ; na parte de Conhecimento Específico, a média foi 62,4.
A maior média desse componente foi obtida pelos estudantes da região Sul ( 69,2 ), e a menor, pelos da região Norte ( 57,5 ).
Quanto à variabilidade das notas, o desvio-padrão de todo o Brasil foi 24,5.
O maior desvio-padrão foi encontrado na regão Centro-Oeste ( 26,0 ), e o menor, na região Norte ( 23,2 ).
A nota máxima obtida foi 100,0 e a nota mínima foi 0,0 em todas as regiões do país.
No Gráfico 5.9, apresenta -se a distribuição das notas na questão discursiva no componente de Conhecimento Específico.
À luz dos dados do Gráfico 5.9, verifica -se que a moda da distribuição das notas da questão discursiva do componente de Conhecimento Específico ocorre no intervalo ( 50 ; 60 ], com 23,7 \% do total de participantes.
Observa -se um valor considerável no intervalo [ 0 ; 10 ], cerca de 3,8 \% dos participantes, entretanto, vale ressaltar que, nesse intervalo, são contabilizadas também as questões deixadas em branco. 5.3.2.2 COMENTÁRIOS SOBRE AS RESPOSTAS À QUESTÃO DISCURSIVA A questão discursiva apresentada para a área de Zootecnia foi composta por um texto introdutório, seguido de três itens ( a, b e c ) relacionadas ao tema central, indicando que a banca procurava examinar conhecimentos específicos dos candidatos.
O texto motivador enfatizava as práticas de biosseguridade relacionadas ao manejo sanitário de aves de corte e postura, com foco, em um primeiro momento, na gripe aviária.
Ficou bastante claro o objetivo de tratar sobre biosseguridade aplicada à influenza aviária na avicultura comercial brasileira.
A questão abordou um tema atual e de grande importância na avicultura comercial brasileira.
Cerca de 15 \% do frango produzido no mundo é brasileiro, ratificando a importância da avicultura comercial para o mercado externo.
Além disso, a questão foi pertinente, apresentou linguagem acessível e abordou um conteúdo basilar da área de Zootecnia, sendo considerada como uma questão de dificuldade média.
A maioria das respostas observadas tratavam sobre o conteúdo solicitado, indicando que o conteúdo exigido era de amplo conhecimento no meio acadêmico.
O item a ) da questão pedia a citação de quatro práticas de biosseguridade para a prevenção da gripe aviária, exigindo dos candidatos conhecimento de práticas de biosseguridade além das tradicionalmente aplicadas na avicultura comercial.
O padrão de respostas indicava uma lista de 24 ações como exemplo.
O item b ) pedia a explicação da conduta imediata a ser tomada em caso de constatação de animal com suspeita de gripe aviária.
Para esse item, o padrão de resposta exigia que o estudante indicasse que a conduta imediata a ser tomada, em caso de constatação de uma ave com suspeita da doença, era notificar aos órgãos de Defesa Sanitária Animal competentes.
No item c ) deveria ser apontado o principal prejuízo à avicultura brasileira caso fosse identificado algum surto de gripe aviária em criações comerciais do país.
O padrão de resposta ( pontuação integral ) deveria focar nas barreiras comerciais ou de exportação, ou seja, nos impedimentos de venda de produtos avícolas em caso de ocorrência da doença nos plantéis comerciais.
Considerando o padrão de resposta estabelecido para o item a ), em geral, as respostas dos estudantes apresentaram certa concordância com o padrão esperado.
No entanto, algumas delas não atenderam ao que foi solicitado, seja por estarem incompletas, não mencionando os quatro itens essenciais de biosseguridade para prevenção da influenza aviária, seja por se limitarem a termos genéricos, a práticas que não estão relacionadas com a biosseguridade ou a práticas aplicáveis somente após a constatação de casos de influenza aviária em granja.
A ocorrência dessas respostas pode estar relacionadas à falta de familiaridade dos estudantes com as práticas específicas e adequadas a serem adotadas nesse contexto.
De maneira geral, o desempenho dos estudantes foi considerado mediano.
No que diz respeito ao item a ), muitos estudantes responderam parcialmente à questão, incluindo itens que não estavam diretamente relacionados à biosseguridade.
Alguns também mencionaram medidas a serem tomadas em caso de detecção positiva de influenza aviária na granja, o que indica uma falta de compreensão do enunciado.
Quanto ao item b ), os estudantes apresentaram práticas de biosseguridade apropriadas para lidar com a influenza aviária na granja, mas demonstraram desconhecimento sobre qual órgão deve ser notificado em caso de detecção da doença, apesar de reconhecerem que é uma obrigação.
Já em relação ao item c ), algumas respostas foram parciais, indicando falta de conhecimento sobre os prejuízos Além disso, o comando solicitava apenas quatro exemplos, mas alguns respondentes escreveram um maior número, assim como os explicaram, aparentemente, buscando garantir maior chance de acerto da questão.
Em relação ao item b ), os estudantes apresentaram respostas de acordo com o padrão estabelecido.
No entanto, muitos estudantes indicaram órgãos que não devem ser acionados em casos de detecção de influenza aviária, tais como Ministério Público, Vigilância Sanitária, Emater, Anvisa, Ministério da Saúde e da Fazenda.
Contudo, a falta de compreensão do conceito de biosseguridade pode ter prejudicado a resposta do item b ), que solicitava que os candidatos indicassem a conduta imediata diante da constatação de uma ave com suspeita de influenza aviária.
Outra hipótese é a de que tenham ocorrido dificuldades de entender se a resposta deveria contemplar as ações dentro da granja ( de manejo ) ou os protocolos sanitários oficiais ( notificações ).
Quanto ao item c ), a maioria das respostas atendeu à expectativa da resposta padrão estabelecida, ou seja, indicaram as barreiras comerciais e/ou de exportação como principal prejuízo à avicultura nacional em caso de surto de Gripe Aviária.
Por outro lado, as respostas também evidenciaram os efeitos diretos no mercado interno ( econômicos, sociais, de segurança alimentar, entre outros ), sem citar as barreiras comerciais e/ou de exportação, divergindo da resposta padrão estabelecida, mas recebendo nota parcial conforme previsto nos critérios de pontuação.
Muitos limitaram suas respostas apenas à perda de animais no lote e, de forma genérica, a “ prejuízos econômicos ”.
Por outro lado, várias respostas indicaram, de forma bastante pertinente e sensata, outros impactos negativos importantes, que não apenas ou principalmente as barreiras comerciais ou a limitação nas exportações em caso de surto oficial de gripe aviária.
Foram citados aspectos como : a alta mortalidade ou necessidade de execução dos animais infectados, prejuízos diretos aos produtores rurais, desabastecimento do mercado interno de carnes e ovos, alteração da oferta/demanda e aumento de preços destes produtos, insegurança alimentar da população brasileira e mundial e marketing negativo à avicultura brasileira.
De maneira geral, o desempenho dos estudantes foi considerado mediano.
No que diz respeito ao item a ), muitos estudantes responderam parcialmente à questão, incluindo itens que não estavam diretamente relacionados à biosseguridade.
Alguns também mencionaram medidas a serem tomadas em caso de detecção positiva de influenza aviária na granja, o que indica uma falta de compreensão do enunciado.
Quanto ao item b ), os estudantes apresentaram práticas de biosseguridade apropriadas para lidar com a influenza aviária na granja, mas demonstraram desconhecimento sobre qual órgão deve ser notificado em caso de detecção da doença, apesar de reconhecerem que é uma obrigação.
Já em relação ao item c ), algumas respostas foram parciais, indicando falta de conhecimento sobre os prejuízos relacionados à exportação e/ou às barreiras comerciais decorrentes da influenza aviária.
Isso sugere uma lacuna no entendimento dos impactos econômicos relacionados a essa doença.
Ficou evidente uma dificuldade de interpretação de texto e de produção de escrita legível por parte de alguns respondentes, além da falta de organização das respostas de acordo com a sequência correta estabelecida.
Muitos começaram apresentando a resposta do item c ), seguido pelo item b ) e deixando o item a ) por último.
Além disso, alguns trocaram as respostas, escrevendo o que deveria ser do item a ) no item b ) e vice-versa. 5.3.2.3 CONSIDERAÇÕES FINAIS SOBRE A QUESTÃO DISCURSIVA DO COMPONENTE ESPECÍFICO A análise global do desempenho dos estudantes na questão discursiva da área de Zootecnia revelou um desempenho mediano.
Foram identificados erros de escrita, dificuldade na organização das respostas e problemas de caligrafia, que prejudicaram a compreensão do conteúdo redigido.
Embora a maioria dos estudantes tenha conseguido compreender e responder às questões solicitadas, algumas respostas estavam incompletas.
Outros demonstraram ter um domínio limitado sobre o assunto, enquanto alguns responderam e justificaram suas respostas.
Foram cerca de 2,5 \% de provas em branco ( sem respostas ), indicando falta de conhecimentos específicos sobre o que foi questionado ou desinteresse no processo \textbf{avaliativo}.
Alguns estudantes também se desviaram do tema da avicultura, discorrendo, por exemplo, sobre política no Brasil.
Outros limitaram -se a comentar sobre o nível de dificuldade da questão, sem abordar os conteúdos específicos solicitados.
Essas falhas demonstram a falta de capacidade por parte de alguns estudantes de estruturar e apresentar as respostas de forma clara e organizada.
Já em relação à apreciação da questão, é importante ressaltar que a questão abordou conteúdos pertinentes para a formação do profissional zootecnista.

% matched lemmas: avaliativo, interseção, sequenciado
\end{textsample}
